%------------------------------------------------------------------------------%
\begin{exercises}

% Exercício 27 página 26
%------------------------------------------------------------------------------%
\begin{exercise}
A série converge ou diverge? Justifique sua resposta.
\[
  \sum_{n=1}^{\infty} \frac{\sqrt{n}}{\;\ln n\;}
\]
\begin{answer}
  Aplicando o teste da divergência temos
  \begin{align*}
  	\lim_{n\to\infty} \frac{\sqrt{n}}{\;\ln n\;} & = \lim_{x\to\infty} \frac{\; x^{\frac{1}{2}} \;}{ \ln x }                \\[2mm]
  	                                             & = \lim_{x\to\infty} \dfrac{\; \frac{1}{2} x^{\frac{-1}{2}} \;}{ x^{-1} } \\[2mm]
  	                                             & = \lim_{x\to\infty} \frac{\; \sqrt{x} \;}{ 2 } = \infty
  \end{align*}
  Como o termo geral não vai para zero a série diverge.
\end{answer}
\end{exercise}

% Exercício 28 página 20
%------------------------------------------------------------------------------%
\begin{exercise}
  Use o teste do $n$-ésimo termo para divergência para mostrar que a série é divergente
  ou afirmar que o teste não é conclusivo.
  \[
  \sum_{n=1}^{\infty} \frac{n(n+1)}{\;(n+2)(n+3)\;}
  \]
  \begin{answer}
    Calculando o limite do dos termos da série temos
    \begin{align*}
    \lim_{n\to\infty} \frac{n(n+1)}{\;(n+2)(n+3)\;} & = \lim \frac{n^2+n}{\;n^2 +5n +6\;} \\
    & = \lim \frac{2n+1}{\;2n +5\;}       \\
    & = \lim \frac{2}{\;2\;} = 1
    \end{align*}
    em cada passo foi identificado que o limite gerava uma
    indeterminação do tipo $\sfrac{\infty}{\infty}$
    e usado o método de L'Hôpital.
    Como o limite do termo geral não é zero a série diverge.
  \end{answer}
\end{exercise}

% 49-68 página 20 do Thomas Vol. 2
%------------------------------------------------------------------------------%
\begin{exercise}
  Verifique se a série é convergente, caso seja calcule seu limite.
  \begin{mathlist}
    \item \sum_{n=0}^{\infty} \left(\frac{1}{\sqrt{2}}\right)^n
    \item \sum_{n=0}^{\infty} \left(\sqrt{2}\right)^n
    \item \sum_{n=1}^{\infty} (-1)^{n+1}\frac{3}{2^n}
    \item \sum_{n=1}^{\infty} (-1)^{n+1}n
    \item \sum_{n=0}^{\infty} \cos(n\pi)
    \item \sum_{n=0}^{\infty} \frac{\cos(n\pi)}{5^n}
    \item \sum_{n=0}^{\infty} e^{-2n}
    \item \sum_{n=1}^{\infty} \ln\left(\frac{1}{3^n}\right)
    \item \sum_{n=1}^{\infty} \frac{2}{10^n}
    \item \sum_{n=0}^{\infty} \frac{1}{x^n} \quad |x|>1
    \item \sum_{n=0}^{\infty} \frac{2^n-1}{3^n}
    \item \sum_{n=1}^{\infty} \left(1-\frac{1}{n}\right)^n
    \item \sum_{n=0}^{\infty} \frac{n!}{1000^n}
    \item \sum_{n=1}^{\infty} \frac{n^n}{n!}
    \item \sum_{n=1}^{\infty} \frac{2^n+3^n}{4^n}
    \item \sum_{n=1}^{\infty} \frac{2^n+4^n}{3^n+4^n}
    \item \sum_{n=1}^{\infty} \ln\left(\frac{n}{n+1}\right)
    \item \sum_{n=1}^{\infty} \ln\left(\frac{n}{2n+1}\right)
    \item \sum_{n=0}^{\infty} \left(\frac{e}{\pi}\right)^n
    \item \sum_{n=0}^{\infty} \frac{e^{n\pi}}{\pi^{ne}}
  \end{mathlist}
  \begin{answer}
    \begin{alphalist}[3]
      \item $2+\sqrt{2}$
      \item Diverge
      \item $1$
      \item Diverge
      \item Diverge
      \item $\sfrac{5}{6}$
      \item $\sfrac{e^2}{e^2-1}$
      \item Diverge
      \item $\sfrac{2}{9}$
      \item $\sfrac{x}{x-1}$
      \item $\sfrac{3}{2}$
      \item Diverge
      \item Diverge
      \item Diverge
      \item $4$
      \item Diverge
      \item Diverge
      \item Diverge
      \item $\sfrac{\pi}{\pi-e}$
      \item Diverge
    \end{alphalist}
  \end{answer}
\end{exercise}


%------------------------------------------------------------------------------%
%\begin{exercise}
%\begin{answer}
%\end{answer}
%\end{exercise}

%------------------------------------------------------------------------------%
%\begin{exercise}
%\begin{mathlist}[2]
%  \item
%  \item
%  \item
%\end{mathlist}
%\begin{answer}
%  \begin{alphalist}[3]
%    \item
%    \item
%    \item
%  \end{alphalist}
%\end{answer}
%\end{exercise}

\end{exercises}
%------------------------------------------------------------------------------%
